\documentclass[11pt, letterpaper]{article}

\usepackage[margin=1in]{geometry}
\usepackage{amsmath, amssymb, amsthm}
\usepackage{booktabs}
\usepackage{hyperref}
\usepackage{float}
\usepackage{enumitem}
\usepackage{xcolor}
\usepackage{setspace}
\usepackage{titlesec}
\usepackage{parskip}
\usepackage{tcolorbox}
\usepackage{mdframed}
\usepackage{array}

\hypersetup{colorlinks=true, linkcolor=blue!60!black, urlcolor=blue!60!black}

\titlespacing*{\section}{0pt}{1.4em}{0.6em}
\titlespacing*{\subsection}{0pt}{1.0em}{0.4em}

\setstretch{1.15}

% ── Custom boxes ──────────────────────────────────────────────────────────────
\tcbuselibrary{skins, breakable}

\newtcolorbox{whybox}{
  colback=blue!5!white, colframe=blue!50!black,
  fonttitle=\bfseries, title=Why we used this method,
  breakable, left=6pt, right=6pt, top=4pt, bottom=4pt
}

\newtcolorbox{mathbox}{
  colback=gray!5!white, colframe=gray!50!black,
  fonttitle=\bfseries, title=Mathematical formulation,
  breakable, left=6pt, right=6pt, top=4pt, bottom=4pt
}

\newtcolorbox{resultbox}{
  colback=green!5!white, colframe=green!50!black,
  fonttitle=\bfseries, title=Our result,
  breakable, left=6pt, right=6pt, top=4pt, bottom=4pt
}

\newtcolorbox{speakerbox}{
  colback=yellow!8!white, colframe=orange!70!black,
  fonttitle=\bfseries, title=Speaker note,
  breakable, left=6pt, right=6pt, top=4pt, bottom=4pt
}

% ── Theorem environments ───────────────────────────────────────────────────────
\theoremstyle{definition}
\newtheorem{theorem}{Theorem}[section]
\newtheorem{lemma}[theorem]{Lemma}
\newtheorem{corollary}[theorem]{Corollary}
\newtheorem{definition}[theorem]{Definition}

% ──────────────────────────────────────────────────────────────────────────────

\title{\textbf{Presentation Scripts \& Method Justifications}\\[0.4em]
\large Does Negative Forum Sentiment Predict Declines in Active Player Counts?\\[0.3em]
\normalsize MATH 189: Data Analysis \& Inference, Winter 2026}
\author{Jiho Lee, Kevin Xu, Viki Shi, Sharon Tey, Favio Espejo, Ankita Inamti\\
\small University of California, San Diego}
\date{}

\begin{document}
\maketitle
\thispagestyle{empty}

\begin{speakerbox}
\textbf{How to use this document.}
Each section covers one method from the report.
Every section has four parts:
(1) a plain-language description of what the method does,
(2) why we chose it for this project,
(3) the key mathematical formulation, and
(4) what we found.
The speaker notes give suggested wording for the presentation.
\end{speakerbox}

\tableofcontents
\newpage

% ══════════════════════════════════════════════════════════════════════════════
\section{Research Question and Study Design}

\textbf{What we are asking:} Does negative user sentiment on Steam forums predict short-term declines in how many people are actively playing a game?

\textbf{Data:} 15 Steam games, January 2020 -- September 2024 (weekly observations). Games are stratified into three lifecycle groups: Popular (CS2, Dota 2, PUBG, GTA V, Warframe), Decline (Rust, ARK, For Honor, Dead by Daylight, Rainbow Six Siege), and Volatile (No Man's Sky, Destiny 2, Terraria, Monster Hunter: World, PAYDAY 2).

\begin{whybox}
We stratified the sample deliberately to avoid survivorship bias.
If we only picked the most popular games on Steam right now, we would miss declining games entirely, and our results would only apply to blockbusters.
By including games from all three lifecycle phases we get variation in both sentiment and player counts that is necessary for statistical power.
\end{whybox}

\begin{speakerbox}
``We started by asking a simple question: when players on Steam complain a lot, does that actually drive people away?
To answer it fairly, we needed a diverse sample --- not just the biggest games.
So we used stratified sampling: five popular titles, five in decline, and five volatile ones with big player swings.
This gives us a complete picture of the gaming ecosystem.''
\end{speakerbox}

% ══════════════════════════════════════════════════════════════════════════════
\section{Sentiment Construction: RoBERTa + CLT}

\textbf{What it is:} Each Steam review is classified by a pre-trained RoBERTa transformer model (\texttt{cardiffnlp/twitter-roberta-base-sentiment}) into three classes: Positive, Neutral, Negative. The model outputs a probability score $s_j \in [0,1]$ for the ``Negative'' class. Reviews are aggregated into a weekly score per game.

\begin{whybox}
We used RoBERTa instead of a simpler rule-based approach (like VADER) for two reasons.
First, RoBERTa uses \emph{contextual embeddings}: it reads words in context, so sarcasm, gaming slang (``this patch is trash''), and negation (``not bad'') are handled more accurately.
Second, a pre-trained transformer requires no labeled training data from us --- we just run it on the Steam reviews directly.

The weekly aggregation lets us use the \textbf{Central Limit Theorem} to justify treating $S_{it}$ as approximately normal for inference, even though individual scores are skewed.
\end{whybox}

\begin{mathbox}
Let $\mathcal{P}_{it}$ be the set of reviews for game $i$ in week $t$.
The weekly negative sentiment score is the mean of individual scores:
\[
    S_{it} = \frac{1}{|\mathcal{P}_{it}|} \sum_{j \in \mathcal{P}_{it}} s_j \in [0,1]
\]

\begin{theorem}[Lindeberg--L\'evy CLT]
Let $s_1, s_2, \ldots, s_n \overset{\text{i.i.d.}}{\sim} F$ with $\mathbb{E}[s_j]=\mu$ and $\mathrm{Var}(s_j)=\sigma^2 < \infty$.
Then:
\[
    \sqrt{n}\,\frac{\bar{s}_n - \mu}{\sigma} \xrightarrow{d} \mathcal{N}(0,1) \quad \text{as } n \to \infty
\]
\end{theorem}

\textbf{Implication:} Even though individual review scores are skewed, $S_{it}$ is approximately normal for large $|\mathcal{P}_{it}|$.
This justifies standard $t$-tests, $F$-tests, and confidence intervals applied to $S_{it}$.
\end{mathbox}

\begin{speakerbox}
``To measure sentiment, we ran every Steam review through a state-of-the-art language model called RoBERTa.
It reads each review in context --- so it understands sarcasm and gaming slang --- and outputs a negativity score between 0 and 1.
We then average those scores for each game each week.
The Central Limit Theorem tells us that even though individual scores are messy, the weekly average behaves like a normal distribution --- and that's what makes all our downstream statistics valid.''
\end{speakerbox}

% ══════════════════════════════════════════════════════════════════════════════
\section{Stationarity Testing: Augmented Dickey--Fuller}

\textbf{What it is:} A statistical test for whether a time series has a \emph{unit root} (non-stationary trend) or is stationary (fluctuates around a stable mean). Non-stationary series produce spurious correlations in regression.

\begin{whybox}
Ordinary regression assumes that the variables are stationary.
If we regressed non-stationary player counts directly on non-stationary sentiment, we could find a ``significant'' relationship purely because both series happen to trend downward over time --- even if they have nothing to do with each other.
This is called \emph{spurious regression}.

We tested all 30 series (2 variables $\times$ 15 games) with the ADF test, then first-differenced the player count series to remove the trend before regression.
\end{whybox}

\begin{mathbox}
The ADF test fits the model:
\[
    \Delta y_t = \alpha + \delta\, y_{t-1} + \sum_{j=1}^{p} \phi_j\, \Delta y_{t-j} + \varepsilon_t
\]

$H_0: \delta = 0$ (unit root, non-stationary) \quad vs \quad $H_1: \delta < 0$ (stationary).

\textbf{Key fact:} Under $H_0$, the test statistic does \emph{not} follow a normal or $t$-distribution.
It follows the Dickey--Fuller distribution, which is left-skewed with more negative critical values (5\% critical value $\approx -2.86$ with intercept, vs.\ $-1.645$ for normal).
\end{mathbox}

\begin{resultbox}
25 of 30 series have unit roots at $\alpha = 0.05$.
5 stationary series: Rust (\texttt{log\_players}) and ARK, Dead by Daylight, Monster Hunter: World, Warframe (\texttt{neg\_sentiment}).
Solution: we use $\Delta\log(\text{players})$ (first-differenced) as the dependent variable throughout.
\end{resultbox}

\begin{speakerbox}
``Before running any regressions, we had to check whether our time series are stationary --- meaning they fluctuate around a stable mean rather than drifting up or down.
If we skip this step and both series happen to trend downward, we'd find a fake correlation.
The Augmented Dickey--Fuller test checks for this.
It turns out 25 of our 30 series are non-stationary, so we take first differences of the log player counts.
That removes the trend and gives us a clean dependent variable: the week-over-week change.''
\end{speakerbox}

% ══════════════════════════════════════════════════════════════════════════════
\section{Granger Causality Tests}

\textbf{What it is:} A test of whether past values of sentiment help predict future player changes, \emph{beyond} what past player changes alone can predict. It tests for \emph{temporal precedence}, not structural causation.

\begin{whybox}
We cannot run a controlled experiment on Steam games.
Granger causality is the standard tool for establishing predictive precedence in observational time series data.
A significant Granger test tells us that sentiment contains forward-looking information about player behavior --- even after accounting for the game's own momentum.

We ran separate Granger tests for each of the 15 games, which lets us see which specific titles drive the effect.
\end{whybox}

\begin{mathbox}
\begin{definition}[Granger Causality]
$\{S_t\}$ Granger-causes $\{Y_t\}$ if lagged $S$ improves the mean-squared forecast of $Y$ beyond lagged $Y$ alone.
\end{definition}

\textbf{Restricted model} (own lags only):
\[
    Y_t = c + \sum_{j=1}^{p} \phi_j\, Y_{t-j} + \varepsilon_t^R
\]

\textbf{Unrestricted model} (own lags + sentiment lags):
\[
    Y_t = c + \sum_{j=1}^{p} \phi_j\, Y_{t-j} + \sum_{j=1}^{q} \psi_j\, S_{t-j} + \varepsilon_t^U
\]

$H_0: \psi_1 = \psi_2 = \cdots = \psi_q = 0$.

\begin{theorem}[Granger F-Test]
Under $H_0$ with i.i.d.\ normal errors:
\[
    F = \frac{(\mathrm{RSS}_R - \mathrm{RSS}_U)/q}{\mathrm{RSS}_U/(T - 2p - q - 1)} \sim F(q,\; T - 2p - q - 1)
\]
\end{theorem}

We test at lags 1--4 for each game.
\end{mathbox}

\begin{resultbox}
\textbf{7 of 15 games} reject $H_0$ at $\alpha = 0.05$.
Strongest: For Honor ($F = 39.197$, $p < 0.001$) and PUBG ($F = 19.823$, $p < 0.001$).
An 8th game (Counter-Strike 2) is marginally significant at $\alpha = 0.10$.
\end{resultbox}

\begin{speakerbox}
``The Granger causality test answers: does knowing last week's sentiment give us any extra information about this week's player count, beyond what last week's player count already tells us?
We run this separately for all 15 games and find that 7 of them show significant results.
This tells us sentiment isn't just a side effect --- it actually carries information about where player counts are heading.''
\end{speakerbox}

% ══════════════════════════════════════════════════════════════════════════════
\section{Permutation Test for Granger Causality}

\textbf{What it is:} A non-parametric version of the Granger test that does not assume normally distributed errors.
It generates a reference distribution by randomly shuffling the sentiment time series and recomputing the $F$-statistic.

\begin{whybox}
The parametric Granger $F$-test assumes i.i.d.\ normal errors, which our diagnostics show is violated (Shapiro--Wilk $W = 0.923$, $p < 0.001$; Ljung--Box $p < 0.001$).
The permutation test is \emph{distribution-free}: its $p$-value is exact in finite samples regardless of the error distribution.
If a game passes the permutation test, the result cannot be blamed on a distributional assumption.
\end{whybox}

\begin{mathbox}
\textbf{Procedure} ($B = 2{,}000$ replications):
\begin{enumerate}[nosep]
    \item Randomly permute $\{S_{it}\}_{t=1}^T$ within game $i$, destroying temporal order.
    \item Recompute the Granger $F$-statistic on the permuted data: $F^{*(b)}$.
    \item Repeat $B$ times.
\end{enumerate}

\[
    p_{\text{perm}} = \frac{1}{B}\sum_{b=1}^{B} \mathbf{1}\!\left(F^{*(b)} \geq F_{\text{obs}}\right)
\]

\begin{lemma}[Exact Validity]
$p_{\text{perm}}$ is exact under $H_0$ in finite samples.
This follows from the exchangeability of the permuted series under the null of no temporal dependence between $S$ and $Y$.
\end{lemma}
\end{mathbox}

\begin{resultbox}
\textbf{4 of 15 games} survive at $\alpha = 0.05$: For Honor, PUBG, GTA V, Rainbow Six Siege.
Two others are marginal: Counter-Strike 2 ($p_{\text{perm}} = 0.075$) and Monster Hunter: World ($p_{\text{perm}} = 0.063$).
\end{resultbox}

\begin{speakerbox}
``We followed up the Granger test with a permutation test, which is more robust because it makes no distributional assumptions.
We randomly shuffle the sentiment time series --- scrambling the temporal link between sentiment and player counts --- and recompute the $F$-statistic 2,000 times.
The question is: how often does the shuffled data produce an $F$ as large as what we actually observed?
Only 4 games pass this stricter test, which tells us the strongest results are real, not statistical artifacts.''
\end{speakerbox}

% ══════════════════════════════════════════════════════════════════════════════
\section{Two-Way Fixed Effects Panel Regression}

\textbf{What it is:} A regression model that controls for all time-invariant differences between games (e.g., genre, studio quality) and for all time-specific shocks that affect all games simultaneously (e.g., Steam seasonal sales). This is the main model of the paper.

\begin{whybox}
We have panel data: 15 games observed over 243 weeks.
A simple OLS regression on the pooled data would confound the sentiment effect with, say, the fact that some games are inherently more popular.
Fixed effects remove both types of confounding:
\begin{itemize}[nosep]
    \item \textbf{Game fixed effects} ($\alpha_i$): absorb all time-invariant game characteristics (genre, studio reputation, initial player base).
    \item \textbf{Time fixed effects} ($\gamma_t$): absorb platform-wide events that affect all games at the same time (Steam Summer Sale, COVID lockdowns, competing game releases).
\end{itemize}
After removing these effects, the remaining variation in $\hat{\beta}$ reflects only the within-game relationship between sentiment and player changes.
\end{whybox}

\begin{mathbox}
The model:
\[
    \Delta\log Y_{it} = \alpha_i + \gamma_t + \beta\, S_{i,t-1} + X_{it}\,\theta + \varepsilon_{it}
\]

\textbf{Within (demeaning) estimator.} Define:
\[
    \ddot{Z}_{it} = Z_{it} - \bar{Z}_{i\cdot} - \bar{Z}_{\cdot t} + \bar{Z}_{\cdot\cdot}
\]
OLS on demeaned data eliminates $\alpha_i$ and $\gamma_t$:
\[
    \hat{\beta}_{\text{FE}} = \left(\sum_{i,t} \ddot{\mathbf{x}}_{it}\ddot{\mathbf{x}}_{it}'\right)^{-1} \sum_{i,t} \ddot{\mathbf{x}}_{it}\ddot{Y}_{it}
\]

\begin{corollary}[Consistency \& Asymptotic Normality]
Under strict exogeneity $\mathbb{E}[\varepsilon_{it} \mid \mathbf{x}_i, \alpha_i, \gamma_t] = 0$:
\[
    \sqrt{N}\,(\hat{\beta}_{\text{FE}} - \beta) \xrightarrow{d} \mathcal{N}\!\left(0,\; \Sigma^{-1}\Omega\,\Sigma^{-1}\right)
\]
where $\Omega$ accounts for within-game serial correlation (clustered standard errors).
\end{corollary}
\end{mathbox}

\begin{resultbox}
Three models estimated ($N = 3{,}630$ observations):
\begin{center}
\small
\begin{tabular}{lccc}
\toprule
& \textbf{Model 1} & \textbf{Model 2} & \textbf{Model 3} \\
& Entity FE & Two-Way FE & Multi-Lag \\
\midrule
$\hat{\beta}_1$ & $-0.511^{**}$ $(0.251)$ & $-0.649^{**}$ $(0.316)$ & $-0.152$ $(0.151)$ \\
$\hat{\beta}_2$ & --- & --- & $-0.017$ $(0.024)$ \\
$\hat{\beta}_3$ & --- & --- & $-0.611^{**}$ $(0.268)$ \\
$p$-value & $0.042$ & $0.040$ & --- \\
\bottomrule
\end{tabular}
\end{center}

\textbf{Model 2 preferred:} $\hat{\beta} = -0.649$, $p = 0.040$, 95\% CI $= [-1.268,\,-0.030]$.\\
\textbf{Model 3 cumulative:} $-0.152 - 0.017 - 0.611 = -0.780$ over three weeks.
\end{resultbox}

\begin{speakerbox}
``The core of our analysis is a two-way fixed effects regression.
Think of it this way: we strip out everything that makes each game unique --- its genre, its fanbase, its history --- and we also strip out everything that happened on Steam platform-wide in a given week.
What's left is the pure within-game, within-week relationship between sentiment and player changes.
Our estimate: a one-unit increase in negativity is associated with a 64.9\% decline in that week's player count change.
The $p$-value is 0.040, which is significant, and the confidence interval excludes zero.''
\end{speakerbox}

% ══════════════════════════════════════════════════════════════════════════════
\section{Hausman Test: Fixed vs.\ Random Effects}

\textbf{What it is:} A specification test that determines whether game-level unobservables are correlated with the regressors (fixed effects model required) or uncorrelated (random effects is also valid).

\begin{whybox}
The random effects (RE) model is more efficient when its assumptions hold, because it uses between-game variation in addition to within-game variation.
However, RE assumes that unobserved game characteristics (e.g., community toxicity, developer responsiveness) are \emph{uncorrelated} with sentiment --- a strong and implausible assumption.
We ran the Hausman test to check this formally.
Even though the test fails to reject the RE assumptions, we still prefer FE on theoretical grounds.
\end{whybox}

\begin{mathbox}
\begin{theorem}[Hausman, 1978]
Under $H_0$ (RE correctly specified, $\text{Cov}(\alpha_i, \mathbf{x}_{it}) = 0$):
\[
    H = (\hat{\boldsymbol{\beta}}_{\text{FE}} - \hat{\boldsymbol{\beta}}_{\text{RE}})'\!
    \Big[\widehat{\text{Var}}(\hat{\boldsymbol{\beta}}_{\text{FE}}) - \widehat{\text{Var}}(\hat{\boldsymbol{\beta}}_{\text{RE}})\Big]^{-1}\!
    (\hat{\boldsymbol{\beta}}_{\text{FE}} - \hat{\boldsymbol{\beta}}_{\text{RE}}) \xrightarrow{d} \chi^2_k
\]
Reject $H_0$ $\Rightarrow$ only FE is consistent. Fail to reject $\Rightarrow$ RE is also acceptable.
\end{theorem}
\end{mathbox}

\begin{resultbox}
$\chi^2(2) = 1.774$, $p = 0.412$: fail to reject $H_0$.
Random effects is statistically acceptable, but we prefer FE because the RE assumption is implausible in this context --- game-specific unobservables (community size, developer behavior) are plausibly correlated with sentiment levels.
\end{resultbox}

\begin{speakerbox}
``We asked: do we really need fixed effects, or could a simpler random effects model work?
The Hausman test compares the two: if game-specific unobservables correlate with sentiment, random effects gives biased estimates and only fixed effects is safe.
The test $p$-value of 0.412 says statistically we can't tell the difference --- but theoretically, we know game characteristics like studio reputation are correlated with how players review the game.
So we stuck with fixed effects as the safer, more defensible choice.''
\end{speakerbox}

% ══════════════════════════════════════════════════════════════════════════════
\section{Clustered Standard Errors}

\textbf{What it is:} A method of computing standard errors that is robust to both heteroscedasticity (non-constant error variance) and serial correlation (errors within a game correlated over time).

\begin{whybox}
Standard OLS standard errors assume errors are i.i.d.\ (independent and identically distributed).
Two diagnostic tests show this fails badly in our data:
\begin{itemize}[nosep]
    \item \textbf{Breusch--Pagan test}: LM $= 30.537$, $p < 0.001$ --- heteroscedasticity present.
    \item \textbf{Ljung--Box test}: $Q(20) = 31{,}353$, $p < 0.001$ --- strong serial correlation.
\end{itemize}
Using ordinary SEs would make our confidence intervals too narrow and our $t$-statistics too large, leading to over-rejection of $H_0$.
Clustered SEs correct for both problems by allowing any within-game correlation structure.
\end{whybox}

\begin{mathbox}
\begin{theorem}[Clustered Variance Estimator]
When observations are independent across games $i = 1,\ldots,N$ but arbitrarily correlated within games:
\[
    \hat{V}_{\text{clust}} = (X'X)^{-1} \!\left(\sum_{i=1}^{N} \mathbf{X}_i'\hat{\boldsymbol{\varepsilon}}_i\hat{\boldsymbol{\varepsilon}}_i'\mathbf{X}_i\right)\!(X'X)^{-1}
\]
This is consistent for $\text{Var}(\hat{\boldsymbol{\beta}})$ as $N \to \infty$, regardless of within-cluster correlation.
\end{theorem}

\textbf{Breusch--Pagan LM test:} Under $H_0$ (homoscedasticity), regress $\hat{\varepsilon}^2$ on regressors; $\text{LM} = n \cdot R^2 \sim \chi^2_k$.

\textbf{Ljung--Box portmanteau test:}
\[
    Q(m) = T(T+2)\sum_{k=1}^{m} \frac{\hat{\rho}_k^2}{T-k} \xrightarrow{d} \chi^2_m \quad \text{under no serial correlation}
\]
\end{mathbox}

\begin{speakerbox}
``Our standard errors use a technique called clustering.
The idea is that observations from the same game are correlated with each other over time --- if GTA V had a bad week, the next week is probably related.
Standard errors that ignore this would be artificially small, making our results look more significant than they are.
Clustered standard errors allow any correlation pattern within a game and give us an honest picture of uncertainty.''
\end{speakerbox}

% ══════════════════════════════════════════════════════════════════════════════
\section{Block Bootstrap Confidence Intervals}

\textbf{What it is:} A resampling method that constructs confidence intervals by repeatedly re-estimating the model on random resamples of the data, without assuming any parametric distribution for the errors.

\begin{whybox}
Asymptotic confidence intervals rely on the normal approximation $\hat{\beta} \approx \mathcal{N}(\beta, \text{SE}^2)$.
With $N = 15$ games and non-normal, serially correlated errors, this approximation may be poor.
The block bootstrap:
\begin{itemize}[nosep]
    \item Requires \emph{no distributional assumptions} on $\varepsilon_{it}$.
    \item Resamples entire game panels (not individual rows) to preserve within-game time structure.
    \item Provides an independent check: if the bootstrap CI agrees with the asymptotic CI, both are trustworthy.
\end{itemize}
\end{whybox}

\begin{mathbox}
\textbf{Procedure} ($B = 2{,}000$ replications, seed $= 189$):
\begin{enumerate}[nosep]
    \item Draw 15 games with replacement from $\{1,\ldots,15\}$.
    \item Stack all time periods for the sampled games into a bootstrap panel.
    \item Fit the entity-FE model; record $\hat{\beta}^{*(b)}$.
    \item Repeat. Report $[\hat{\beta}^{*}_{(0.025)},\, \hat{\beta}^{*}_{(0.975)}]$.
\end{enumerate}

\begin{lemma}[Why game-level resampling?]
Resampling entire game panels preserves the within-game temporal dependence $\{\varepsilon_{it}\}_{t=1}^{T}$.
Row-level resampling would treat observations as i.i.d.\ and understate uncertainty when serial correlation is present.
\end{lemma}

Bootstrap SE: $\displaystyle\widehat{\text{SE}}_{\text{boot}} = \sqrt{\frac{1}{B-1}\sum_{b=1}^{B}\!\left(\hat{\beta}^{*(b)} - \bar{\hat{\beta}}^*\right)^2}$
\end{mathbox}

\begin{resultbox}
Mean $\hat{\beta}^* = -0.512$, bootstrap SE $= 0.268$.\\
95\% bootstrap CI $= [-1.093,\,-0.050]$ --- excludes zero.\\
95\% asymptotic CI $= [-1.004,\,-0.019]$ --- also excludes zero.\\
Both CIs agree: the negative effect is robust to distributional assumptions.
\end{resultbox}

\begin{speakerbox}
``To make sure our confidence interval doesn't depend on assuming normal errors, we also computed a bootstrap CI.
We resample all 15 games with replacement 2,000 times, re-estimate the model each time, and look at the distribution of estimates.
The 2.5th and 97.5th percentiles give us the CI.
The key design choice is resampling whole games --- not individual weeks --- so we preserve the time structure within each game.
Our bootstrap CI of $[-1.093, -0.050]$ agrees with the asymptotic one, confirming the result is real.''
\end{speakerbox}

% ══════════════════════════════════════════════════════════════════════════════
\section{Statistical Power Analysis}

\textbf{What it is:} Power is the probability of correctly rejecting $H_0$ when it is false. The Minimum Detectable Effect (MDE) is the smallest true effect size for which the test achieves 80\% power.

\begin{whybox}
We report power because a statistically significant result with low power still deserves scrutiny.
With only $N = 15$ games, our study is limited in its ability to detect small effects.
Reporting power is \emph{honest}: it tells the reader how much to trust a null result, and it tells the reader that our significant result ($p = 0.040$) is at the edge of what our design can reliably detect.
\end{whybox}

\begin{mathbox}
\begin{theorem}[Power Formula]
For a two-sided test at level $\alpha$ with $\hat{\beta} \approx \mathcal{N}(\beta_a, \text{SE}^2)$:
\[
    \text{Power}(\beta_a) = 1 - \Phi\!\left(z_{1-\alpha/2} - \frac{|\beta_a|}{\text{SE}}\right) + \Phi\!\left(-z_{1-\alpha/2} - \frac{|\beta_a|}{\text{SE}}\right)
\]
\end{theorem}

\textbf{Minimum Detectable Effect.} The smallest $|\beta^*|$ achieving power $\geq 1-\kappa$:
\[
    |\beta^*| = (z_{1-\alpha/2} + z_{1-\kappa})\cdot\text{SE}(\hat{\beta})
\]

\textbf{Derivation:} Set Power $= 1-\kappa$ (ignoring negligible second tail):
\[
    z_{1-\alpha/2} - \frac{|\beta^*|}{\text{SE}} = -z_{1-\kappa}
    \implies |\beta^*| = (z_{1-\alpha/2} + z_{1-\kappa})\cdot\text{SE}
\]
\end{mathbox}

\begin{resultbox}
With SE $= 0.251$, $\alpha = 0.05$ ($z_{0.975} = 1.960$), 80\% power ($z_{0.80} = 0.842$):
\[
    \text{MDE} = (1.960 + 0.842)\times 0.251 = 0.707
\]
Our estimated effect $|\hat{\beta}| = 0.511 < 0.707 = \text{MDE}$, so power at the observed effect is only $\mathbf{0.535}$.
The design is \textbf{underpowered}: even at the true effect, we would correctly reject $H_0$ only about half the time.
\end{resultbox}

\begin{speakerbox}
``Here's where we have to be honest about our study's limitations.
Power is the probability that we would detect a real effect if it exists.
With 15 games, our power is only 0.535 at the observed effect size --- barely above a coin flip.
The minimum detectable effect is 0.707, but our estimate is 0.511, which is smaller.
So while we did find a significant result, we need many more games in the panel to be confident we're not missing anything.
We report this transparently so the reader can judge the result appropriately.''
\end{speakerbox}

% ══════════════════════════════════════════════════════════════════════════════
\section{Stratum-Specific Regressions}

\textbf{What it is:} Re-estimating Model 2 separately for each of the three game lifecycle strata (Popular, Decline, Volatile) to check whether the sentiment--player relationship varies by game type.

\begin{whybox}
The panel regression estimates a single $\hat{\beta}$ that is an average across all 15 games.
But the effect might be very different for a blockbuster like CS2 vs.\ a struggling game like For Honor.
Stratum-specific regressions reveal this heterogeneity and tell us \emph{for which types of games} sentiment is most predictive.

This also serves as a robustness check: if the negative sign only shows up in one stratum, the overall result is less convincing.
\end{whybox}

\begin{resultbox}
\begin{center}
\begin{tabular}{lcccc}
\toprule
\textbf{Stratum} & $\hat{\beta}$ & SE & $p$-value & Interpretation \\
\midrule
Popular  & $+0.022$ & 0.151 & 0.884 & Near-zero; large games are resilient \\
Decline  & $-0.595$ & 0.454 & 0.189 & Negative but not significant (5 games) \\
Volatile & $-1.369^{**}$ & 0.570 & 0.016 & Strongest effect; fragile player bases \\
\bottomrule
\end{tabular}
\end{center}
Sentiment is most predictive for volatile games where the player base is already fragile and responsive to community discourse.
\end{resultbox}

\begin{speakerbox}
``When we split the analysis by game type, an interesting pattern emerges.
For popular games --- CS2, Dota 2, GTA V --- the estimate is essentially zero.
These games have huge, loyal player bases that don't respond much to a few bad reviews.
For volatile games --- the ones with big swings around updates --- sentiment has a much stronger effect: an estimate of negative 1.37.
This makes intuitive sense: in a volatile game, the marginal player is more likely to be swayed by community sentiment.''
\end{speakerbox}

% ══════════════════════════════════════════════════════════════════════════════
\section{LOESS: Nonparametric Regression Check}

\textbf{What it is:} Locally Estimated Scatterplot Smoothing (LOESS) fits a smooth curve through a scatterplot without assuming any global functional form.

\begin{whybox}
Our panel model assumes a \emph{linear} relationship between sentiment and player changes.
If the true relationship is curved (e.g., sentiment only matters above a threshold, or the effect reverses at very high negativity), the linear model would be misspecified and the slope estimate $\hat{\beta}$ would be meaningless.

LOESS lets us visually verify linearity without imposing it.
If the LOESS curve tracks the OLS line, linear regression is appropriate.
\end{whybox}

\begin{mathbox}
At each evaluation point $x_0$, LOESS fits a local weighted polynomial:
\[
    \hat{m}(x_0) = \arg\min_{\mathbf{b}} \sum_{i=1}^{n} K\!\left(\frac{|x_i - x_0|}{h(x_0)}\right)\!\left(y_i - b_0 - b_1(x_i - x_0)\right)^2
\]
where $K(u) = (1 - |u|^3)^3\,\mathbf{1}(|u| \leq 1)$ is the tricube kernel and $h(x_0)$ is the adaptive bandwidth (span $= 0.3$).

No global functional form is assumed; the curve is derived entirely from local patterns in the data.
\end{mathbox}

\begin{resultbox}
The LOESS curve (span $= 0.3$) closely tracks the OLS line across all 3,630 observations. No evidence of nonlinearity. The linear specification is validated.
\end{resultbox}

\begin{speakerbox}
``One assumption we had to verify was linearity.
LOESS is a nonparametric smoother --- it fits a curve through the scatterplot without any assumption about its shape.
If the true relationship were curved or threshold-like, the LOESS line would deviate from our OLS line.
But they track each other closely, which tells us the linear model is appropriate for this data.''
\end{speakerbox}

% ══════════════════════════════════════════════════════════════════════════════
\section{Placebo Test: Reversed Lag Direction}

\textbf{What it is:} A falsification test that checks for reverse causation by testing whether \emph{future} sentiment predicts \emph{past} player changes.

\begin{whybox}
A concern with any observational study is reverse causation: maybe declining player counts \emph{cause} negative sentiment, not the other way around.
If this is the only mechanism, then future sentiment should also ``predict'' past changes (because both are driven by the same underlying decline).
By reversing the lag direction, we test this.

A genuine temporal precedence result should show many more significant games in the original direction than in the reversed direction.
\end{whybox}

\begin{resultbox}
\begin{itemize}[nosep]
    \item Original direction (sentiment leads changes): \textbf{7/15} games significant at $\alpha = 0.05$.
    \item Reversed direction (future sentiment predicts past changes): \textbf{3/15} games significant.
\end{itemize}
The asymmetry (7 vs.\ 3) provides partial evidence against pure reverse causation.
However, 3 significant reversed results caution that some bidirectional feedback may exist.
\end{resultbox}

\begin{speakerbox}
``A natural concern is whether we have the direction backwards --- maybe declining player counts cause negative reviews, not the other way around.
To test this, we flip the lag direction: does \emph{next} week's sentiment predict \emph{this} week's player change?
Under pure reverse causation, both directions should give similar results.
But we get 7 significant games going forward and only 3 going backward.
That asymmetry supports our interpretation that sentiment has genuine predictive content for future player behavior, not just as a reflection of the past.''
\end{speakerbox}

% ══════════════════════════════════════════════════════════════════════════════
\section{Out-of-Sample Prediction}

\textbf{What it is:} Splitting the panel into a training period and a test period, then evaluating how well the model predicts player changes on data it has never seen.

\begin{whybox}
In-sample $R^2$ and $p$-values can be inflated by overfitting.
Out-of-sample prediction tests whether sentiment provides genuine \emph{practical} predictive value beyond a naïve baseline.
The naïve baseline here is predicting zero change every week (i.e., using only the mean).

A small out-of-sample improvement is honest evidence about effect magnitude.
\end{whybox}

\begin{resultbox}
\begin{center}
\begin{tabular}{lccc}
\toprule
\textbf{Metric} & \textbf{Model} & \textbf{Naïve Baseline} & \textbf{Improvement} \\
\midrule
RMSE & 0.509 & 0.510 & 0.2\% \\
MAE  & 0.391 & 0.392 & 0.2\% \\
\bottomrule
\end{tabular}
\end{center}
Training: 2020-02-10 to 2023-05-29 (2,587 obs).
Test: 2023-06-05 to 2024-09-30 (1,043 obs).

The 0.2\% improvement is statistically detectable but practically weak.
Sentiment is a real signal but a minor one compared to other factors driving player behavior.
\end{resultbox}

\begin{speakerbox}
``The final practical test: does our model actually predict better than doing nothing?
We trained on data through May 2023 and predicted the held-out period through September 2024.
The sentiment model beats a naïve baseline by 0.2\% in RMSE.
That's a real improvement, but it's tiny.
Sentiment is a statistically detectable signal, but it explains only a small part of what drives player behavior.
Game updates, competitor launches, streaming trends --- all of these matter much more.''
\end{speakerbox}

% ══════════════════════════════════════════════════════════════════════════════
\section{Summary: Why Each Method?}

\begin{center}
\renewcommand{\arraystretch}{1.35}
\small
\begin{tabular}{p{3.8cm} p{4.2cm} p{5.5cm}}
\toprule
\textbf{Method} & \textbf{Why we used it} & \textbf{What it told us} \\
\midrule
Stratified sampling & Avoid survivorship bias; ensure variation in outcome & Captured popular, declining, and volatile games fairly \\
\midrule
RoBERTa sentiment & Contextual understanding of gaming language; no labeled data needed & Outperforms bag-of-words models on nuanced review text \\
\midrule
CLT justification & Makes weekly mean $S_{it}$ approximately normal for inference & Validates use of $t$- and $F$-tests on aggregated sentiment \\
\midrule
ADF stationarity & Prevents spurious regression from trending series & 25/30 series non-stationary; first-differencing required \\
\midrule
Granger causality & Tests temporal precedence in observational time series & 7/15 games show sentiment leads player changes \\
\midrule
Permutation test & Distribution-free alternative to parametric Granger & 4/15 games survive; strongest results are not artifacts \\
\midrule
Two-way FE & Controls for game and time heterogeneity simultaneously & $\hat{\beta} = -0.649$, $p = 0.040$; robust to confounding \\
\midrule
Hausman test & Formally compare FE vs.\ RE specifications & FE preferred on both statistical and theoretical grounds \\
\midrule
Clustered SEs & Correct for heteroscedasticity and serial correlation & Honest uncertainty quantification for panel time series \\
\midrule
Block bootstrap & Non-parametric CI; no distributional assumption & CI $= [-1.093, -0.050]$; confirms asymptotic result \\
\midrule
Power analysis & Honest reporting of study limitations & Power $= 0.535$; MDE $= 0.707$; design underpowered \\
\midrule
Stratum regressions & Test effect heterogeneity across game lifecycle phases & Effect concentrated in volatile games ($\hat{\beta} = -1.369$) \\
\midrule
LOESS & Verify linearity assumption nonparametrically & Linear specification validated \\
\midrule
Placebo test & Falsification against reverse causation & 7 vs.\ 3 significant games; temporal ordering supported \\
\midrule
Out-of-sample prediction & Test practical predictive value; guard against overfitting & 0.2\% improvement; signal is real but weak \\
\bottomrule
\end{tabular}
\end{center}

\vspace{1em}
\begin{speakerbox}
\textbf{Closing script suggestion:}\\
``Every method we used was chosen for a specific reason.
We didn't just throw statistical tests at the data --- each one addresses a specific threat to the validity of our conclusions.
The Granger test checks temporal ordering.
The permutation test drops the normality assumption.
Fixed effects remove game-level confounding.
The bootstrap gives us distribution-free intervals.
And the power analysis tells us honestly what we can and cannot conclude with 15 games.
The overall picture is consistent: negative sentiment has a statistically significant but practically modest predictive effect on player counts, and it's strongest for games where the player base is already fragile.''
\end{speakerbox}

\end{document}
